\section{A smile is a probability law}
\label{sec:intro}

A listed option market gives finitely many prices.  A smile model must turn
them into a complete curve: a price and an implied volatility at every strike,
including strikes that were not quoted.  The completed curve will be
differentiated for digitals and densities, integrated for variance contracts,
and compared with curves at other expiries.  It is therefore not enough to
draw a smooth line through the quotes.

The obstruction is elementary.  The second strike derivative of a call price
is a probability density.  A curve that fits every quote can nevertheless
make this derivative negative between two quotes.  It then assigns a negative
price to a butterfly and a negative mass to an interval of terminal prices.
No numerical routine has malfunctioned; the curve was allowed to leave the
space of probability laws.

The useful question is thus not ``which curve fits well?'' but ``in which
coordinates is every fitted curve already a law?''  The answer developed here
uses quantiles.  In rank coordinates, unit mass is automatic and positive
density reduces to monotonicity.  Taking the logarithm of the quantile speed
removes even that constraint.  The resulting log-quantile-density (LQD)
family provides unconstrained coordinates for a broad class of mean-one laws
with exponential log-return tails.

\subsection{Normalized market}

Fix one expiry $T$.  Let $F$ be its forward and $D$ its discount factor.  Under
the $T$-forward measure define
\begin{equation}
X=\log\frac{S_T}{F},\qquad Y=e^X=\frac{S_T}{F},\qquad \E[Y]=1.
\label{eq:Xdef}
\end{equation}
The last equality is the martingale condition.  Write
$k=\log(K/F)$ for log-moneyness.  Normalized call and put prices are
\begin{equation}
c(k)=\E[(Y-e^k)^+],\qquad P(k)=\E[(e^k-Y)^+],
\qquad c(k)-P(k)=1-e^k.
\label{eq:calldef}
\end{equation}
A dollar price is $DF$ times its normalized price.

Let $\tau$ denote variance time: calendar time, possibly adjusted for the
desk's event-time convention.  Total implied variance is
$w(k)=\sigma(k)^2\tau$.  The normalized Black call is
\begin{equation}
\Black(k,w)=\PhiN(d_+)-e^k\PhiN(d_-),
\qquad d_\pm=-\frac{k}{\sqrt w}\pm\frac{\sqrt w}{2},
\label{eq:black}
\end{equation}
and $w(k)$ is the root of $\Black(k,w(k))=c(k)$
\citep{Black1976}.  Rates and carry have now disappeared from the single-slice
problem.

\subsection{The static constraint set}

Use normalized strike $y=e^k$.  If $Y$ has density $f_Y$, differentiation of
the payoff gives
\begin{equation}
\frac{\partial c}{\partial y}=-\Prob(Y>y),
\qquad
\frac{\partial^2c}{\partial y^2}=f_Y(y)\ge0.
\label{eq:BL}
\end{equation}
This is the Breeden--Litzenberger identity
\citep{BreedenLitzenberger1978}.  It identifies the geometry of the call curve
with the law of $Y$: decreasing slope is a digital probability and convexity
is positive density.

The identity can be read using traded portfolios.  For a strike spacing
$\delta$, the symmetric butterfly
\[
c(y-\delta)-2c(y)+c(y+\delta)
\]
is the price of a long call below, a long call above, and two short calls at
the center.  Dividing by $\delta^2$ and shrinking the spacing recovers
$f_Y(y)$.  A negative second difference is therefore not merely undesirable
curvature; it is a negative price for a nonnegative terminal payoff.  The
probability statement and the static-arbitrage statement are the same.

\begin{definition}[Arbitrage-free slice]
\label{def:slice}
An arbitrage-free slice is a call curve of the form \eqref{eq:calldef} for a
positive random variable $Y$ with $\E[Y]=1$.  When $Y$ has a density, this is
equivalent to a decreasing convex function of strike $y$ satisfying
\[
(1-y)^+\le c(y)\le1,\qquad -1\le c'(y)\le0,\qquad
c(y)\longrightarrow0\quad(y\to\infty),
\]
with $c'(y)\to-1$ as $y\downarrow0$ when there is no mass at zero.
\end{definition}

Each boundary condition has a probabilistic source.  The upper bound follows
from $(Y-y)^+\le Y$ and $\E[Y]=1$.  The lower bound is Jensen's inequality,
$(\E[Y]-y)^+\le\E[(Y-y)^+]$.  The slope is a survival probability and must
lie between $-1$ and zero.  Finally, calls vanish at infinite strike because
the integrable random variable $Y$ carries no mean mass arbitrarily far in
its tail.  A valid call curve is thus exactly a compact encoding of a
mean-one law.

Convexity belongs to $y$, not to $k$.  Indeed,
\[
\frac{\partial^2c}{\partial k^2}
=e^k\{e^kf_Y(e^k)-\Prob(Y>e^k)\},
\]
which is negative far in the left wing for any ordinary density.  A convexity
test performed in log-strike will therefore reject correct call curves.

The same constraint becomes less transparent in volatility coordinates.
Differentiating $c(k)=\Black(k,w(k))$ twice yields
\begin{equation}
f_X(k)=g_{\rm D}(k)\frac{\phin(d_-(k,w(k)))}{\sqrt{w(k)}},
\quad
g_{\rm D}(k)=
\left(1-\frac{k w'}{2w}\right)^2
-\frac{(w')^2}{4}\left(\frac1w+\frac14\right)+\frac{w''}{2}.
\label{eq:durrleman}
\end{equation}
Thus a volatility curve is butterfly-free precisely when
$g_{\rm D}(k)\ge0$ for all $k$ \citep{GatheralJacquier2014}.  The condition is
global, nonlinear, and involves two derivatives.  Direct volatility models
must enforce it; the representation below makes it an identity.

This explains why good residuals alone are weak evidence.  Quotes constrain
only finitely many values of $w$.  The function $g_{\rm D}$ also depends on
the slope and curvature between them, and may become negative in a narrow
unquoted interval.  Increasing the strike grid reduces but does not remove
the logical problem: the optimizer still moves in a space containing invalid
curves.  The quantile construction reverses the order of work.  It first
chooses a law and only then asks which smile the Black formula assigns to it.

Two checks fix the normalization in \eqref{eq:durrleman}.  If the smile is
flat, $w'=w''=0$, so $g_{\rm D}\equiv1$ and the remaining factor is exactly
the Black density of $X$.  At the money, the identities derived later in
\cref{sec:ledger} give
$g_{\rm D}(0)=\sqrt{w_0}f_X(0)/\phin(\sqrt{w_0}/2)$.  Thus
$g_{\rm D}\ge0$ is not a penalty formula: it is positive density written in
the volatility chart.

\subsection{Notation ledger}

The construction uses three horizontal axes.  Strike $y$ is the axis on which
calls are convex; log-moneyness $k=\log y$ is the market axis; rank $u$ and
log-odds $z$ are the model axes.  For a rank $u\in(0,1)$, the \emph{odds} are
$u/(1-u)$ and the \emph{log-odds} are
\[
z=\log\frac{u}{1-u}.
\]
The function taking rank to log-odds is called the \emph{logit}:
$\logit(u)=\log\{u/(1-u)\}$.  It maps $(0,1)$ bijectively onto $\R$; its
inverse is the logistic function $\Lambda(z)=1/(1+e^{-z})$.  Thus $u$ and $z$
record the same percentile on two different scales.  The recurring notation
is collected below.

\begin{table}[htbp]
\centering\small
\setlength{\tabcolsep}{4pt}
\begin{tabular}{@{}llll@{}}
\toprule
$X$, $Y=e^X$ & log and gross forward return &
$u$, $z=\logit u$ & rank and log-odds\\
$k$, $y=e^k$ & log and normalized strike &
$\Lambda$, $\rho$ & logistic CDF and density\\
$c$, $P$ & normalized call and put &
$Q$, $q=Q'$ & quantile and quantile density\\
$w=\sigma^2\tau$ & total implied variance &
$g$ & smooth log-speed in rank\\
$x$, $m$ & transport and martingale shift &
$G$ & upper-share ledger\\
$L,R,a_n$ & polynomial coordinates &
$\lambda_-,\lambda_+$ & endpoint speeds\\
$u_k,z_k$ & strike rank and log-odds &
$u_*,f_*,\Delta$ & ATM rank, density, digital gap\\
$\beta_L,\beta_R$ & Lee wing slopes &
$N,Z_{\max}$ & basis order and grid half-width\\
\bottomrule
\end{tabular}
\caption{Notation used throughout the chapter.  A prime denotes an ordinary
derivative of a one-variable function; $\partial$ denotes a partial
derivative.  The time $\tau$ is always variance time.}
\label{tab:notation}
\end{table}

\subsection{What the chapter separates}

Three questions recur, and they should not be confused.  \emph{Validity} asks
whether a parameter vector defines a probability law.  \emph{Information}
asks which features of that law are determined by the quoted strip.
\emph{Convention} selects the remainder: basis order, regularization, tail
coordinates, and the scope of calendar control.  The LQD coordinates settle
the first question structurally.  The later sections measure the second and
state the third explicitly.

The development is progressive.  \Cref{sec:model} constructs the coordinates.
\Cref{sec:validity,sec:tails} establish their probabilistic content.
\Cref{sec:ledger} derives prices and local smile quantities from one integral.
\Cref{sec:calibration,sec:calendar} turn the representation into a fitter and a
surface.  \Cref{sec:examples} then shows what the construction says on market
and synthetic data.  Proofs and implementation details are collected in the
appendices.
